% ============================================================
\chapter{Span, Independence, Bases, and Dimension}
\label{ch:span-bases-dimension}
% ============================================================

\section*{Chapter Preview}
\addcontentsline{toc}{section}{Chapter Preview}

A small list of vectors can describe a large space. Span explains how. Linear
independence explains when no vector in the list is redundant. A basis is a
list that is just large enough to describe every vector uniquely.

\section{Linear Combinations and Span}

\begin{definition}
  A \emph{linear combination} of vectors $v_1,\ldots,v_k$ is a vector of the
  form
  \[
    \lambda_1v_1+\cdots+\lambda_kv_k,
  \]
  where the coefficients $\lambda_i$ are scalars.
\end{definition}

\begin{definition}
  The \emph{span} of $v_1,\ldots,v_k$ is the set of all their linear
  combinations:
  \[
    \Span\{v_1,\ldots,v_k\}
    =
    \{\lambda_1v_1+\cdots+\lambda_kv_k:\lambda_i\in\F\}.
  \]
\end{definition}

\begin{proposition}
  $\Span\{v_1,\ldots,v_k\}$ is the smallest subspace containing
  $v_1,\ldots,v_k$.
\end{proposition}

\begin{proof}
  The span is closed under linear combinations, so it is a subspace. It
  contains each $v_i$ by taking coefficient $1$ on $v_i$ and $0$ on the
  others. Any subspace containing all the $v_i$ must contain all their linear
  combinations, so it must contain the span.
\end{proof}

\section{Linear Independence}

\begin{definition}
  A list $v_1,\ldots,v_k$ is \emph{linearly independent} if
  \[
    \lambda_1v_1+\cdots+\lambda_kv_k=0
  \]
  implies $\lambda_1=\cdots=\lambda_k=0$.
\end{definition}

If a list is not independent, then one vector is redundant: it can be written
as a linear combination of the others.

\section{Bases}

\begin{definition}
  A \emph{basis} of a vector space $V$ is a list of vectors that spans $V$ and
  is linearly independent.
\end{definition}

The two conditions work together. Spanning says every vector can be described.
Independence says the description is unique.

\begin{theorem}[Unique coordinates]
  If $b_1,\ldots,b_n$ is a basis of $V$, then every $v\in V$ can be written in
  a unique way as
  \[
    v=c_1b_1+\cdots+c_nb_n.
  \]
\end{theorem}

\begin{proof}
  Existence comes from spanning. For uniqueness, suppose
  \[
    v=c_1b_1+\cdots+c_nb_n=d_1b_1+\cdots+d_nb_n.
  \]
  Subtracting gives
  \[
    0=(c_1-d_1)b_1+\cdots+(c_n-d_n)b_n.
  \]
  Independence forces $c_i-d_i=0$ for all $i$, so $c_i=d_i$.
\end{proof}

\section{Dimension}

\begin{definition}
  If $V$ has a basis with $n$ vectors, then $V$ is \emph{finite-dimensional}
  and $\dim V=n$.
\end{definition}

The fact that every basis of the same finite-dimensional space has the same
number of vectors is a theorem. It is what makes dimension well-defined.

\section{Python: Coordinates in a Basis}

If the basis vectors are the columns of a matrix $B$, then finding coordinates
of $v$ means solving
\[
  Bc=v.
\]

\begin{lstlisting}
import numpy as np

B = np.array([[1.0, 1.0],
              [1.0, -1.0]])
v = np.array([3.0, 1.0])

c = np.linalg.solve(B, v)
print("coordinates:", c)
print("reconstruction:", B @ c)
\end{lstlisting}

\begin{computebox}
The command \texttt{np.linalg.solve(B, v)} calls optimized numerical linear
algebra routines. Later we will learn the mathematical structure behind this
kind of computation.
\end{computebox}

\section*{Chapter Summary}
\addcontentsline{toc}{section}{Chapter Summary}

\begin{itemize}
  \item Span is the set of all linear combinations.
  \item Linear independence means no nontrivial linear combination gives zero.
  \item A basis is a spanning independent list.
  \item A basis gives unique coordinates.
  \item Dimension is the number of vectors in any basis.
\end{itemize}

\section*{Where This Appears Later}
\addcontentsline{toc}{section}{Where This Appears Later}

Coordinates and bases are the mechanism behind matrices, change of basis,
diagonalization, Fourier methods, and PCA.

\section*{Exercises}
\addcontentsline{toc}{section}{Exercises}

\subsection*{A. Theory}

\begin{exercise}
\exerciseid{4.A1}
Prove that the span of a list of vectors is a subspace.
\end{exercise}

\begin{exercise}
\exerciseid{4.A2}
Prove that if a list of vectors is linearly dependent, then at least one vector
in the list is a linear combination of the others.
\end{exercise}

\begin{exercise}
\exerciseid{4.A3}
Explain why a basis gives unique coordinates.
\end{exercise}

\subsection*{B. By Hand}

\begin{exercise}
\exerciseid{4.B1}
Do the vectors $(1,0,1)$, $(0,1,1)$, and $(1,1,2)$ form a basis of $\R^3$?
Explain without doing unnecessary arithmetic.
\end{exercise}

\begin{exercise}
\exerciseid{4.B2}
Let $b_1=(1,1)$ and $b_2=(1,-1)$. Find the coordinates of $v=(3,1)$ in the
basis $(b_1,b_2)$.
\end{exercise}

\begin{exercise}
\exerciseid{4.B3}
Find a basis for the span of
\[
  (1,2,0),\quad (2,4,0),\quad (0,1,1).
\]
\end{exercise}

\subsection*{C. Python / NumPy}

\begin{exercise}
\exerciseid{4.C1}
Put the vectors in Exercise 4.B3 into the columns of a matrix and use
\texttt{np.linalg.matrix\_rank} to check the dimension of their span.
\end{exercise}

\begin{exercise}
\exerciseid{4.C2}
Write a function \texttt{coords(B, v)} that returns the coordinate vector
$c$ solving $Bc=v$.
\end{exercise}

\subsection*{D. Challenge}

\begin{exercise}
\exerciseid{4.D1}
Give two different bases of $\mathcal P_2$ and compute the coordinates of
$p(x)=1+x+x^2$ in both bases.
\end{exercise}

